\section*{Chapter \ref{ch:qm:finite-difference-methods} --- Finite-Difference Methods}

\begin{solution}{exo:qm:finite-difference-methods:1}
$\Delta\tau \le \Delta x^2/\sigma^2 = 6.25 \times 10^{-6}/0.04 = 1.5625 \times 10^{-4}$ years, so 6\,400 steps for one year.
\end{solution}

\begin{solution}{exo:qm:finite-difference-methods:2}
$4.3572 + (4.3572 - 4.3558)/3 = 4.35767$.
\end{solution}

\begin{solution}{exo:qm:finite-difference-methods:3}
$u(x \pm h) = u \pm hu_x + \frac{h^2}{2}u_{xx} \pm \frac{h^3}{6}u_{xxx} + \frac{h^4}{24}u_{xxxx} + O(h^5)$; adding, subtracting $2u$ and dividing by $h^2$ leaves $u_{xx} + \frac{h^2}{12}u_{xxxx} + O(h^4)$.
\end{solution}

\begin{solution}{exo:qm:finite-difference-methods:4}
$g(\pi) = (1 - 64)/(1 + 64) = -0.969$; $|g|^n = 0.01$ needs $n = \ln100/(-\ln0.969) \approx 148$ steps, far more than the 25 of the whole run. One implicit half step at $\lambda/2 = 16$ has factor $1/(1 + 64) =
0.015$ at $\xi = \pi$: a single half step damps that mode by 65.
\end{solution}

\begin{solution}{exo:qm:finite-difference-methods:5}
The coefficients are $a/h^2 \mp b/(2h)$; both are nonnegative iff $h \le 2a/|b|$, i.e.\ $|b|h/(2a) \le 1$.
\end{solution}

\begin{solution}{exo:qm:finite-difference-methods:6}
Far below the strike the call is worth almost exactly zero, so $u = 0$ is nearly the true value and its small error is damped by diffusion before it reaches the strike. Far above, the call is worth
about $S - Ke^{-r\tau}$, a large number: imposing zero there is a large error that diffuses inward and, over the option's life, reaches the strike.
\end{solution}

\begin{solution}{exo:qm:finite-difference-methods:7}
\texttt{digital\_placement()}: on a node $9.9 \times 10^{-3}$, $4.9 \times 10^{-3}$, $2.5 \times 10^{-3}$, $1.2 \times 10^{-3}$ (ratios 2.0, first order); between nodes $6.6 \times 10^{-6}$, $1.5 \times 10^{-6}$,
$3.8 \times 10^{-7}$, $9.5 \times 10^{-8}$ (ratios 4.4, 4.0, 4.0, second order).
\end{solution}

\begin{solution}{exo:qm:finite-difference-methods:8}
A converged price says nothing about its derivatives: a node-to-node oscillation of size $\epsilon$ in the values becomes $\epsilon/h$ in delta and $\epsilon/h^2$ in gamma. In the chapter's run the price is within
0.018 and the gamma wrong by twenty times its size. Greeks need their own convergence test.
\end{solution}

\begin{solution}{pb:qm:finite-difference-methods:1}
\textbf{1.} $\lambda = a\Delta\tau/\Delta x^2 = 0.02 \times 0.01/0.0025^2 = 32$; the \omterm{def:qm:finite-difference-methods:theta}{explicit scheme} needs $\lambda \le \frac12$.
\textbf{2.} At $\lambda = \frac12$: at 0.489 (409 steps on 200 nodes) the price is 4.3576; at 0.509 (393 steps) the largest value on the grid is 2\,413.
\textbf{3.} $(1 - 2\lambda)/(1 + 2\lambda) = -63/65 = -0.969$.
\textbf{4.} The price is within 0.018 of 4.3576; the gamma swings between $-0.81$ and $0.55$ around 0.040 (largest error 0.84).
\textbf{5.} It doubles: 0.46, 0.84, 1.69 and 3.38 on 200 to 1\,600 nodes, because $\lambda$ doubles.
\textbf{6.} Largest gamma error $4.7 \times 10^{-4}$ at 400 nodes, halving with each refinement: first order.
\textbf{7.} $1.76 \times 10^{-5}$, $4.35 \times 10^{-6}$, $1.07 \times 10^{-6}$: ratios 4.05 and 4.07, second order.
\textbf{8.} First order: 0.041, 0.018, 0.0092, 0.0046.
\textbf{9.} A discontinuity sampled at a node (value $\frac12$ there) leaves an error at the strike proportional to the space step, as the measured ratios of 2 show; the loss comes from how
the discontinuous data are sampled, not from the time stepping, since smaller time steps do not change it.
\textbf{10.} Error $6.6 \times 10^{-6}$ on 200 nodes, falling by 4 per halving.
\textbf{11.} About 2.9 times smaller error at equal node count (0.0079 against 0.0225 on 50 nodes).
\textbf{12.} From $1.8 \times 10^{-3}$ and $4.6 \times 10^{-4}$ to $3.2 \times 10^{-7}$. Without the start-up the price error is first order, and extrapolating with $p = 2$ removes the wrong term.
\textbf{13.} When $|b|\Delta x/(2a) > 1$; here 2.49, and the digital reaches 1.030 (above the discounted payout 0.951) with total variation 1.23; upwinding restores 0.951 and monotonicity.
\textbf{14.} Zero gamma and Dirichlet ends give prices at the strike that differ by $6 \times 10^{-8}$ with the grid at five standard deviations.
\textbf{15.} Errors 0.045, 0.026 and 0.014 against Margrabe's 10.5243 on $40^2$, $80^2$, $160^2$ nodes: first order, from the kink along the diagonal between nodes.
\textbf{16.} A refinement test on the Greeks (gamma should change by a factor 4 less per halving, not grow), or a count of sign changes of the second difference near the strike.
\textbf{17.} Stability bounds the growth of every mode but not its sign: Crank--Nicolson keeps high frequencies alive with alternating sign, and differencing twice amplifies them by $1/h^2$.
\textbf{18.} Implicit start-up on every non-smooth datum, strikes and barriers between nodes, measured convergence orders for prices and Greeks, an \omterm{def:qm:finite-difference-methods:m}{M-matrix} check, and extrapolation only at a measured order.
\textbf{19.} Named result: \emph{the sawtooth gamma}: four implicit half steps remove the oscillation and restore second order, the gamma error falling by 4.05 and 4.07 per halving, where plain Crank--Nicolson's
doubles and two half steps give first order.
\textbf{20.} What the scheme does to each frequency at each step, which is what the Greeks see.
\end{solution}

\begin{solution}{iq:qm:finite-difference-methods:1}
$\Delta t \le \Delta x^2/(2a)$ for $u_t = au_{xx}$ ($\Delta t \le \Delta x^2/\sigma^2$ in log-price). Halving $\Delta x$ quarters the time step, so a fine grid makes the \omterm{def:qm:finite-difference-methods:theta}{explicit scheme} very expensive; hence \omterm{def:qm:finite-difference-methods:theta}{implicit
schemes}.
\iqlookfor{The $\Delta x^2$ scaling and its cost.}
\end{solution}

\begin{solution}{iq:qm:finite-difference-methods:2}
Its amplification factor tends to $-1$ at high frequencies when $\Delta t/\Delta x^2$ is large, so the high-frequency content of a kinked payoff survives with alternating sign and shows up in the second
difference. Fix: a few implicit (half) steps at the start (Rannacher), payoff smoothing or placement, or an L-stable scheme.
\iqlookfor{The amplification factor near $-1$ and the Rannacher start.}
\end{solution}

\begin{solution}{iq:qm:finite-difference-methods:3}
Midway between nodes (or with the payoff averaged over each cell) for strikes; exactly on a node for a continuously monitored barrier, where the boundary condition is imposed. The error of a
discontinuity inside a cell depends on its position in the cell and destroys the convergence order.
\iqlookfor{Placement determines the order.}
\end{solution}

\begin{solution}{iq:qm:finite-difference-methods:4}
In one dimension, the tridiagonal system is solved by the Thomas algorithm in $O(n)$. In two, an ADI splitting solves one tridiagonal system per grid line and per direction, $O(n^2)$ per step, instead of
a sparse solve of the full system.
\iqlookfor{Thomas and ADI.}
\end{solution}

\begin{solution}{iq:qm:finite-difference-methods:5}
Refine space and time together by factors of two on a problem with a known answer (or against a much finer reference), and check that the error ratio tends to $2^p$ for prices and Greeks separately,
at points that stay on the grid.
\iqlookfor{Error ratios per halving, including Greeks.}
\end{solution}

\begin{solution}{iq:qm:finite-difference-methods:6}
The scheme is not monotone: central first differences with a large drift relative to diffusion (or a mixed-derivative stencil with strong correlation) give negative off-diagonal coefficients, and
Crank--Nicolson does not damp the resulting oscillations. Upwind, refine where the drift dominates, use a positivity-preserving mixed-derivative stencil, and start with implicit steps.
\iqlookfor{Loss of the \omterm{def:qm:finite-difference-methods:m}{M-matrix} property.}
\end{solution}
